\documentclass[12pt]{article}
\usepackage[utf8]{inputenc}
\usepackage[T1]{fontenc}
\usepackage[margin=1in]{geometry}
\usepackage{setspace}
\usepackage{xcolor}
\usepackage[colorlinks=true, urlcolor=blue!60!black, linkcolor=blue!60!black]{hyperref}
\setlength{\parindent}{0em}
\setlength{\parskip}{0.6em}
\renewcommand{\baselinestretch}{1.0}
\begin{document}
\begin{center}
{\LARGE \textit{Legal Services Corp. v. Velazquez}, 531 U.S. 533 (2001)\par}
\vspace{0.5em}
{\large Case Brief --- Two-Page Edition\par}
\end{center}

\section{Memory Jogger}
Congress violated the First Amendment by barring federally funded legal aid lawyers from challenging the validity of welfare laws, because the program facilitated private advocacy, not government speech, and distorted the adversary system (\href{https://supreme.justia.com/cases/federal/us/531/533/}{531 U.S. 533}).

\section{Detailed Case Facts}
A 1996 appropriations rider, \S~504(a)(16), barred Legal Services Corporation (LSC) grantees from welfare-reform litigation and lobbying, allowing individual benefit suits only ``if such relief does not involve an effort to amend or otherwise challenge existing law.'' Under LSC rules, grantee lawyers could argue factual or statutory error but not that a welfare law was unconstitutional or preempted, and had to withdraw if such an issue arose mid-case.

\section{Procedural History}
Grantee lawyers, clients, and officials sued in the Eastern District of New York. The Second Circuit struck and severed the suits-for-benefits proviso as viewpoint discrimination; the Supreme Court affirmed.

\section{Judicial Votes}
Affirmed, 5--4. Kennedy wrote for Stevens, Souter, Ginsburg, and Breyer; Scalia dissented, joined by Rehnquist, O'Connor, and Thomas.

\section{Holding}
When a funding program facilitates private speech, Congress may not impose viewpoint-based conditions that distort the bar and the courts by excluding validity arguments and insulating its own laws from review.

\section{Analysis of Opinions}
Kennedy distinguished \textit{Rust v. Sullivan} as a program in which private speakers transmitted the government's message; LSC lawyers instead speak for private clients against the government. The restriction distorted an existing medium of expression, the bar and the courts, creating ``two tiers of cases'' and undermining judicial review, and indigent clients could rarely find substitute counsel. Congress may not relabel a speech restriction a ``definition of its program.'' Scalia, dissenting, saw a symmetrical subsidy condition squarely governed by \textit{Rust}: refusing to fund speech is not abridging it, the ``nondistortion'' principle had no precedent, and severing only the proviso funded litigation Congress had refused to pay for.

\section{Examples: Future Applications}
\textbf{Same-Side:} a state eviction-defense program that bars funded lawyers from challenging housing statutes is invalid. \textbf{Opposite-Side:} a condition barring funded quit-smoking counselors from endorsing vaping conveys the government's own message and is valid. \textbf{Fence-Sitter:} funded asylum clinics barred from systemic challenges but allowed individual constitutional claims.

\section{Critique}
Critics aligned with Scalia say the decision blurs the line between abridging and not subsidizing speech and privileges lawyers over doctors. Access-to-justice advocates praise it for cabining \textit{Rust}, while doctrinal commentators fault its blend of tests, exposed in \textit{American Library Ass'n}.

\section{Key Quotations}
\begin{quote}
``Although the LSC program differs from the program at issue in \textit{Rosenberger} in that its purpose is not to `encourage a diversity of views,' the salient point is that, like the program in \textit{Rosenberger}, the LSC program was designed to facilitate private speech, not to promote a governmental message.''
\end{quote}
\noindent --- Justice Kennedy, Opinion of the Court, 531 U.S. at 542.

\end{document}
